% GENERATED FILE. Edit canonical lessons, then run npm run generate:books.
% canonical-source-hash: fnv1a64:d948677aa9b6eba8
% canonical-lessons: RU-C447-masterskaya, RU-C447-sapozhnik, RU-C447-chasovshchik, RU-C447-slesar, RU-C447-remont

\chapter{Repair shops}
\label{ch:ru-repair-shops}
\hlchaptermodality{447}

\begin{chapteropening}
\textbf{By the end of this chapter:} \emph{I can find repair shops: a workshop, a shoemaker, a watchmaker, a locksmith and repairs.}

Complete the last lesson of chapter 447: I can find repair shops: a workshop, a shoemaker, a watchmaker, a locksmith and repairs.
\end{chapteropening}

\section[\texorpdfstring{\ru{мастерская} (masterskáya)}{мастерская (masterskáya)}]{\ru{мастерская} (masterskáya) — a workshop, a repair shop}

\label{lesson:RU-C447-masterskaya}



\begin{quote}
Before the new one: say the Russian for a drainpipe, then the Russian for a stall.
\end{quote}

\subsection*{You'll want to know: \ru{мастерская}}
\textbf{\ru{мастерская}} (\emph{masterskáya}) — \textquotedblleft{}a workshop, a repair shop\textquotedblright{}.

\begin{quote}
\ru{В} \ru{нашем} \ru{доме} \ru{есть} \ru{маленькая} \ru{мастерская}.
\end{quote}

\begin{quote}
\emph{V náshem dóme yest málenkaya masterskáya.}
\end{quote}

\begin{quote}
There is a small workshop in our building.
\end{quote}

An adjective used as a noun: \emph{\ru{в} \ru{мастерской}}.

\begin{cousinweb}
\emph{\ru{мастерская}} is built on \emph{\ru{мастер}} 'a master, a craftsman' with an adjective ending. EtymDB traces \emph{\ru{мастер}} through Old East Slavic and Old French \emph{maistre}, which goes back to Latin \emph{magister} 'a master'.
\end{cousinweb}

\begin{grammarlens}[title={What you've built}]
One more word for telling, explaining and describing at B1.
\end{grammarlens}

\subsection*{Your turn}
\begin{itemize}
\raggedright
  \item \emph{Say it:} \emph{masterskáya}
  \item \emph{Say it:} \emph{masterskáya}, once more
  \item \emph{Say it:} a sentence of your own with \emph{masterskáya}
  \item \emph{Recall:} say the Russian for a drainpipe, then the Russian for a stall, then say \emph{masterskáya} again
\end{itemize}

\begin{tcolorbox}[breakable,colback=teal!4,colframe=teal!35!black,title={Before you move on}]
What does \ru{мастерская} mean? (\textquotedblleft{}A workshop, a repair shop\textquotedblright{}.) Say it once more.
\end{tcolorbox}
\section[\texorpdfstring{\ru{сапожник} (sapózhnik)}{сапожник (sapózhnik)}]{\ru{сапожник} (sapózhnik) — a shoemaker, a cobbler}

\label{lesson:RU-C447-sapozhnik}



\begin{quote}
Before the new one: say the Russian for a stall, then the Russian for a workshop.
\end{quote}

\subsection*{You'll want to know: \ru{сапожник}}
\textbf{\ru{сапожник}} (\emph{sapózhnik}) — \textquotedblleft{}a shoemaker, a cobbler\textquotedblright{}.

\begin{quote}
\ru{Сапожник} \ru{починил} \ru{мои} \ru{старые} \ru{сапоги}.
\end{quote}

\begin{quote}
\emph{Sapózhnik pochiníl moí stáryye sapogí.}
\end{quote}

\begin{quote}
The cobbler mended my old boots.
\end{quote}

Masculine. Built on \emph{\ru{сапог}} — a boot.

\begin{cousinweb}
\emph{\ru{сапожник}} is built on \emph{\ru{сапог}} 'a boot' with \emph{-\ru{ник}}. \emph{\ru{Сапог}} is inherited from Proto-Slavic \emph{sapogŭ} (EtymDB).
\end{cousinweb}

\begin{grammarlens}[title={What you've built}]
One more word for telling, explaining and describing at B1.
\end{grammarlens}

\subsection*{Your turn}
\begin{itemize}
\raggedright
  \item \emph{Say it:} \emph{sapózhnik}
  \item \emph{Say it:} \emph{sapózhnik}, once more
  \item \emph{Say it:} a sentence of your own with \emph{sapózhnik}
  \item \emph{Recall:} say the Russian for a stall, then the Russian for a workshop, then say \emph{sapózhnik} again
\end{itemize}

\begin{tcolorbox}[breakable,colback=teal!4,colframe=teal!35!black,title={Before you move on}]
What does \ru{сапожник} mean? (\textquotedblleft{}A shoemaker, a cobbler\textquotedblright{}.) Say it once more.
\end{tcolorbox}
\section[\texorpdfstring{\ru{часовщик} (chasovshchík)}{часовщик (chasovshchík)}]{\ru{часовщик} (chasovshchík) — a watchmaker}

\label{lesson:RU-C447-chasovshchik}



\begin{quote}
Before the new one: say the Russian for a workshop, then the Russian for a shoemaker.
\end{quote}

\subsection*{You'll want to know: \ru{часовщик}}
\textbf{\ru{часовщик}} (\emph{chasovshchík}) — \textquotedblleft{}a watchmaker\textquotedblright{}.

\begin{quote}
\ru{Часовщик} \ru{поменял} \ru{батарейку} \ru{в} \ru{часах}.
\end{quote}

\begin{quote}
\emph{Chasovshchík pomenyál bataréyku v chasákh.}
\end{quote}

\begin{quote}
The watchmaker changed the battery in the watch.
\end{quote}

Masculine. Built on \emph{\ru{часы}} — a watch, a clock.

\begin{cousinweb}
\emph{\ru{часовщик}} is built on \emph{\ru{часовой}} 'of clocks', from \emph{\ru{часы}} 'a watch', on \emph{\ru{час}} 'an hour'. EtymDB does not list \emph{\ru{часовщик}}.
\end{cousinweb}

\begin{grammarlens}[title={What you've built}]
One more word for telling, explaining and describing at B1.
\end{grammarlens}

\subsection*{Your turn}
\begin{itemize}
\raggedright
  \item \emph{Say it:} \emph{chasovshchík}
  \item \emph{Say it:} \emph{chasovshchík}, once more
  \item \emph{Say it:} a sentence of your own with \emph{chasovshchík}
  \item \emph{Recall:} say the Russian for a workshop, then the Russian for a shoemaker, then say \emph{chasovshchík} again
\end{itemize}

\begin{tcolorbox}[breakable,colback=teal!4,colframe=teal!35!black,title={Before you move on}]
What does \ru{часовщик} mean? (\textquotedblleft{}A watchmaker\textquotedblright{}.) Say it once more.
\end{tcolorbox}
\section[\texorpdfstring{\ru{слесарь} (slésar)}{слесарь (slésar)}]{\ru{слесарь} (slésar) — a locksmith, a fitter}

\label{lesson:RU-C447-slesar}



\begin{quote}
Before the new one: say the Russian for a shoemaker, then the Russian for a watchmaker.
\end{quote}

\subsection*{You'll want to know: \ru{слесарь}}
\textbf{\ru{слесарь}} (\emph{slésar}) — \textquotedblleft{}a locksmith, a fitter\textquotedblright{}.

\begin{quote}
\ru{Слесарь} \ru{поменял} \ru{замок} \ru{на} \ru{двери}.
\end{quote}

\begin{quote}
\emph{Slésar pomenyál zamók na dverí.}
\end{quote}

\begin{quote}
The locksmith changed the lock on the door.
\end{quote}

Masculine; plural \emph{\ru{слесаря}} or \emph{\ru{слесари}}.

\begin{cousinweb}
\emph{\ru{слесарь}} was borrowed from German \emph{Schlosser} 'a locksmith' (EtymDB), built on \emph{Schloss} 'a lock'.
\end{cousinweb}

\begin{grammarlens}[title={What you've built}]
One more word for telling, explaining and describing at B1.
\end{grammarlens}

\subsection*{Your turn}
\begin{itemize}
\raggedright
  \item \emph{Say it:} \emph{slésar}
  \item \emph{Say it:} \emph{slésar}, once more
  \item \emph{Say it:} a sentence of your own with \emph{slésar}
  \item \emph{Recall:} say the Russian for a shoemaker, then the Russian for a watchmaker, then say \emph{slésar} again
\end{itemize}

\begin{tcolorbox}[breakable,colback=teal!4,colframe=teal!35!black,title={Before you move on}]
What does \ru{слесарь} mean? (\textquotedblleft{}A locksmith, a fitter\textquotedblright{}.) Say it once more.
\end{tcolorbox}
\section[\texorpdfstring{\ru{ремонт} (remónt)}{ремонт (remónt)}]{\ru{ремонт} (remónt) — repairs, a repair; renovation}

\label{lesson:RU-C447-remont}



\begin{quote}
Before the new one: say the Russian for a watchmaker, then the Russian for a locksmith.
\end{quote}

\subsection*{You'll want to know: \ru{ремонт}}
\textbf{\ru{ремонт}} (\emph{remónt}) — \textquotedblleft{}repairs, a repair; renovation\textquotedblright{}.

\begin{quote}
\ru{Ремонт} \ru{туфель} \ru{занял} \ru{два} \ru{дня}.
\end{quote}

\begin{quote}
\emph{Remónt túfel zányal dva dnya.}
\end{quote}

\begin{quote}
Repairing the shoes took two days.
\end{quote}

Masculine. \emph{\ru{Ремонт} \ru{квартиры}} — doing up a flat. The verb is \emph{\ru{ремонтировать}}.

\begin{cousinweb}
\emph{\ru{ремонт}} is usually traced to French \emph{remonte}; EtymDB links it with \emph{remonter} 'to put back up, to restock'. The word probably first named fresh horses for the army.
\end{cousinweb}

\begin{grammarlens}[title={What you've built}]
One more word for telling, explaining and describing at B1.
\end{grammarlens}

\subsection*{Your turn}
\begin{itemize}
\raggedright
  \item \emph{Say it:} \emph{remónt}
  \item \emph{Say it:} \emph{remónt}, once more
  \item \emph{Say it:} a sentence of your own with \emph{remónt}
  \item \emph{Recall:} say the Russian for a watchmaker, then the Russian for a locksmith, then say \emph{remónt} again
\end{itemize}

\begin{tcolorbox}[breakable,colback=teal!4,colframe=teal!35!black,title={Before you move on}]
What does \ru{ремонт} mean? (\textquotedblleft{}Repairs, a repair; renovation\textquotedblright{}.) Say it once more.
\end{tcolorbox}
